% Method section — Label Construction subsection (Part A)
% Paste this into reports/midsem/sections/method.tex, alongside
% Part C's feature-engineering subsection.
%
% Numbers marked [FILL] should be replaced with your actual run's output
% once make_windows.py has been run and its printed summary recorded.

\subsection{Drift Label Construction}
\label{sec:label-construction}

The EuRoC dataset does not provide a pre-computed odometry-drift signal;
it supplies raw IMU measurements and precise ground-truth pose from a
motion-capture system. We construct the regression target — the position
error that a Visual-Inertial Odometry system would accumulate over a
short time window — directly from these two streams.

\subsubsection{Temporal Alignment}
IMU measurements (200\,Hz) and ground-truth pose estimates are recorded
asynchronously. For every IMU sample we attach the temporally nearest
ground-truth sample, discarding IMU samples whose nearest match exceeds
a 20\,ms tolerance. Across the three sequences used in Phase 1
(MH\_01\_easy, MH\_02\_easy, MH\_03\_medium), this discarded between
1.2\% and 2.6\% of IMU samples, concentrated at sequence boundaries
where ground-truth coverage is sparser than IMU coverage.

\subsubsection{Windowing}
The aligned stream is divided into overlapping windows of
$T_w = 1.0\,\mathrm{s}$ (200 IMU samples at 200\,Hz), with a stride of
$T_s = 0.25\,\mathrm{s}$ between consecutive window starts. This yielded
724, 596, and 523 windows for MH\_01\_easy, MH\_02\_easy, and
MH\_03\_medium respectively.

\subsubsection{IMU Strapdown Integration}
For each window, we compute a model-free position estimate by
integrating the raw IMU measurements forward using standard strapdown
inertial mechanization. Starting from the true position $p_0$, velocity
$v_0$, and orientation quaternion $q_0$ at the window's first timestamp
(taken directly from ground truth), we integrate at each IMU sample
$i$ with timestep $\Delta t_i$:
\begin{align}
  a_{\text{world},i} &= R(q_i)\, a_{\text{body},i} + g \\
  v_{i+1} &= v_i + a_{\text{world},i}\, \Delta t_i \\
  p_{i+1} &= p_i + v_i \Delta t_i + \tfrac{1}{2} a_{\text{world},i} \Delta t_i^2 \\
  q_{i+1} &= q_i \otimes \Delta q(\omega_i, \Delta t_i)
\end{align}
where $R(q_i)$ is the rotation matrix corresponding to the current
orientation estimate, $g$ is the world-frame gravity vector, and
$\Delta q(\omega_i, \Delta t_i)$ is the incremental rotation from the
gyroscope reading $\omega_i$.

Two sign/frame conventions — the sign of $g$ and whether $R(q_i)$ maps
body-to-world or world-to-body — are not fixed by the dataset
documentation alone and were determined empirically: we evaluated all
four combinations on a short (0.25\,s) probe window against ground
truth and selected the combination minimizing position error. This
yielded $g = (0,0,-9.81)\,\mathrm{m/s^2}$ under a body-to-world rotation
convention, with a probe-window error of 0.017\,m — one to two orders of
magnitude lower than the three alternative conventions (0.27–0.63\,m),
which we take as strong evidence the correct convention was identified
rather than an arbitrary best-of-four selection.

\subsubsection{Label Definition}
The drift label for a window is the Euclidean distance between the
integrated position estimate $\hat{p}_{\text{end}}$ at the window's final
timestamp and the true position $p_{\text{end}}$ from ground truth:
\begin{equation}
  y_{\text{err}} = \lVert \hat{p}_{\text{end}} - p_{\text{end}} \rVert_2
\end{equation}
The signed per-axis error $(\Delta x, \Delta y, \Delta z)$ is retained
alongside the scalar magnitude for use in the Phase 2 trajectory
correction step and as additional Phase 1 regression targets. Across the
three sequences, mean 1-second-window \emph{raw}-IMU drift was consistent,
ranging from 0.181\,m to 0.186\,m, with maximum single-window error of
0.30\,m on MH\_03\_medium. After oracle bias subtraction the same 1\,s
windows have mean residual 0.018--0.030\,m (MH\_03 still the largest).

The same windows are also integrated after subtracting the dataset's
ground-truth IMU bias estimates. The resulting columns
(\texttt{err\_mag\_m\_bc}, \texttt{err\_dx/dy/dz\_bc}) are an oracle
comparison, not a replacement of the approved raw-IMU target: a deployed
estimator is not handed those biases, and all features remain raw IMU.
Whether the residual is more motion-dependent is measured in the EDA,
not assumed.

\subsubsection{Validation}
We verified the label construction is physically sound by confirming
that mean drift error increases with window length (0.25\,s to 2.0\,s),
consistent with the expected accumulation behavior of open-loop inertial
integration. This check, along with the convention-selection probe
error, is automated in \texttt{tests/test\_label\_sanity.py} and passed
on all sequences prior to any modeling work.
